\chapter{Coordination, the Public, and the Political Economy}
\label{sec:25}
Citizen assemblies, public consultations and ethics boards can make a system better informed, and the assemblies convened on AI have worked best when asked about a thing with a location, an operator and a consequence, such as facial recognition \autocite{adalovelaceinstitute2021biometrics}. None of them can stop a deployment, and read as aggregation instruments they get democracy wrong: a perfectly deliberative process that concludes a demographic should be surveilled, deported or imprisoned hasn't produced a faultless output with a bad result. It has failed as inquiry, because it decided what is done to people outside it. Enforceable rights that don't depend on being popular, a court willing to rule against a majority, and standing for people affected without being enfranchised are what keep the inquiry honest, and they matter most where the boundary of the public excludes the people a decision falls on.

The guarantee and the social-insurance baseline it is enacted with are drafted in Appendix~\ref{sec:A}, with the reason for preferring a guarantee to a transfer. What the floor needs from them is in Chapter~\ref{sec:8}: an exit the employer doesn't price.
